\documentclass{amsart}
% For dealing with references we use the comment environment
\usepackage{verbatim}
\newenvironment{reference}{\comment}{\endcomment}
%\newenvironment{reference}{}{}
\newenvironment{slogan}{\comment}{\endcomment}
\newenvironment{history}{\comment}{\endcomment}

% For commutative diagrams we use Xy-pic
\usepackage[all]{xy}

% We use 2cell for 2-commutative diagrams.
\xyoption{2cell}
\UseAllTwocells

% We use multicol for the list of chapters between chapters
\usepackage{multicol}

% This is generally recommended for better output
\usepackage{lmodern}
\usepackage[T1]{fontenc}

% For cross-file-references

% Package for hypertext links:
\usepackage{hyperref}

% For any local file, say "hello.tex" you want to link to please

% Theorem environments.
%
\theoremstyle{plain}
\newtheorem{theorem}[subsection]{Theorem}
\newtheorem{proposition}[subsection]{Proposition}
\newtheorem{lemma}[subsection]{Lemma}

\theoremstyle{definition}
\newtheorem{definition}[subsection]{Definition}
\newtheorem{example}[subsection]{Example}
\newtheorem{exercise}[subsection]{Exercise}
\newtheorem{situation}[subsection]{Situation}

\theoremstyle{remark}
\newtheorem{remark}[subsection]{Remark}
\newtheorem{remarks}[subsection]{Remarks}

\numberwithin{equation}{subsection}

% Macros
%
\def\lim{\mathop{\mathrm{lim}}\nolimits}
\def\colim{\mathop{\mathrm{colim}}\nolimits}
\def\Spec{\mathop{\mathrm{Spec}}}
\def\Hom{\mathop{\mathrm{Hom}}\nolimits}
\def\Ext{\mathop{\mathrm{Ext}}\nolimits}
\def\SheafHom{\mathop{\mathcal{H}\!\mathit{om}}\nolimits}
\def\SheafExt{\mathop{\mathcal{E}\!\mathit{xt}}\nolimits}
\def\Sch{\mathit{Sch}}
\def\Mor{\mathop{\mathrm{Mor}}\nolimits}
\def\Ob{\mathop{\mathrm{Ob}}\nolimits}
\def\Sh{\mathop{\mathit{Sh}}\nolimits}
\def\NL{\mathop{N\!L}\nolimits}
\def\CH{\mathop{\mathrm{CH}}\nolimits}
\def\proetale{{pro\text{-}\acute{e}tale}}
\def\etale{{\acute{e}tale}}
\def\QCoh{\mathit{QCoh}}
\def\Ker{\mathop{\mathrm{Ker}}}
\def\Im{\mathop{\mathrm{Im}}}
\def\Coker{\mathop{\mathrm{Coker}}}
\def\Coim{\mathop{\mathrm{Coim}}}

% Boxtimes
%
\DeclareMathSymbol{\boxtimes}{\mathbin}{AMSa}{"02}

%
% Macros for moduli stacks/spaces
%
\def\QCohstack{\mathcal{QC}\!\mathit{oh}}
\def\Cohstack{\mathcal{C}\!\mathit{oh}}
\def\Spacesstack{\mathcal{S}\!\mathit{paces}}
\def\Quotfunctor{\mathrm{Quot}}
\def\Hilbfunctor{\mathrm{Hilb}}
\def\Curvesstack{\mathcal{C}\!\mathit{urves}}
\def\Polarizedstack{\mathcal{P}\!\mathit{olarized}}
\def\Complexesstack{\mathcal{C}\!\mathit{omplexes}}
% \Pic is the operator that assigns to X its picard group, usage \Pic(X)
% \Picardstack_{X/B} denotes the Picard stack of X over B
% \Picardfunctor_{X/B} denotes the Picard functor of X over B
\def\Pic{\mathop{\mathrm{Pic}}\nolimits}
\def\Picardstack{\mathcal{P}\!\mathit{ic}}
\def\Picardfunctor{\mathrm{Pic}}
\def\Deformationcategory{\mathcal{D}\!\mathit{ef}}

\setlength{\emergencystretch}{3em}
\begin{document}
\title[Stacks Project, Tag 0A6R]{242 --- Explicit Cech construction of the right adjoint for derived torsion complexes}
\maketitle
\noindent Copyright (C) 2005--2025 Johan de Jong. Stacks Project authors. Source: \href{https://stacks.math.columbia.edu/tag/0A6R}{Tag 0A6R}.

\noindent Editorial difficulty note (not author text). Difficulty H2: The proof constructs the right adjoint on the unbounded derived category by an extended Cech total complex. It proves torsion of its cohomology using a filtered colimit of Koszul complexes and null homotopies, then proves the counit is invertible on the torsion subcategory. These are actual derived-category construction obligations rather than substitution into a local-cohomology formula..

\noindent Editorial provenance note (not author text). History: excerpt from revision \texttt{a0444\allowbreak 6e57e\allowbreak c1fbc\allowbreak 252a8\allowbreak 71afc\allowbreak ec775\allowbreak 2fb28\allowbreak 07b14}, dualizing.tex, lines 1248--1322. Reference presentation was adapted; original-author context and core support blocks were retained or added. The mathematical wording is unchanged. Licensed under GNU FDL 1.2 or later; no invariant sections or cover texts. See ../licenses/COPYING. Context blocks reproduce author definitions and supporting results; they are not additional counted problems.

\noindent
Let $A$ be a ring and let $I \subset A$ be a finitely generated ideal.
Set $Z = V(I) \subset \Spec(A)$. We will construct a functor
\begin{equation}
\label{equation-local-cohomology}
R\Gamma_Z : D(A) \longrightarrow D_{I^\infty\text{-torsion}}(A).
\end{equation}
which is right adjoint to the inclusion functor. For notation
see Section \href{https://stacks.math.columbia.edu/tag/0BJA}{section bad local cohomology (Tag 0BJA)}. The cohomology
modules of $R\Gamma_Z(K)$ are the {\it local cohomology groups
of $K$ with respect to $Z$}.
By Lemma \href{https://stacks.math.columbia.edu/tag/0A6P}{lemma not equal (Tag 0A6P)} this functor will in general {\bf not} be
equal to $R\Gamma_I( - )$ even viewed as functors into $D(A)$.
In Section \href{https://stacks.math.columbia.edu/tag/0BJD}{section local cohomology noetherian (Tag 0BJD)}
we will show that if $A$ is Noetherian, then the two agree.

\begin{lemma}
\label{lemma-local-cohomology-adjoint}
Let $A$ be a ring and let $I \subset A$ be a finitely generated ideal.
There exists a right adjoint $R\Gamma_Z$ (\ref{equation-local-cohomology})
to the inclusion functor $D_{I^\infty\text{-torsion}}(A) \to D(A)$.
In fact, if $I$ is generated by $f_1, \ldots, f_r \in A$, then we have
$$
R\Gamma_Z(K) =
(A \to \prod\nolimits_{i_0} A_{f_{i_0}} \to
\prod\nolimits_{i_0 < i_1} A_{f_{i_0}f_{i_1}}
\to \ldots \to A_{f_1\ldots f_r}) \otimes_A^\mathbf{L} K
$$
functorially in $K \in D(A)$.
\end{lemma}

\begin{proof}
Say $I = (f_1, \ldots, f_r)$ is an ideal.
Let $K^\bullet$ be a complex of $A$-modules.
There is a canonical map of complexes
$$
(A \to \prod\nolimits_{i_0} A_{f_{i_0}} \to
\prod\nolimits_{i_0 < i_1} A_{f_{i_0}f_{i_1}} \to
\ldots \to A_{f_1\ldots f_r}) \longrightarrow A.
$$
from the extended {\v C}ech complex to $A$.
Tensoring with $K^\bullet$, taking associated total complex,
we get a map
$$
\text{Tot}\left(
K^\bullet \otimes_A
(A \to \prod\nolimits_{i_0} A_{f_{i_0}} \to
\prod\nolimits_{i_0 < i_1} A_{f_{i_0}f_{i_1}} \to
\ldots \to A_{f_1\ldots f_r})\right)
\longrightarrow
K^\bullet
$$
in $D(A)$. We claim the cohomology modules of the complex on the left are
$I$-power torsion, i.e., the LHS is an object of
$D_{I^\infty\text{-torsion}}(A)$. Namely, we have
$$
(A \to \prod\nolimits_{i_0} A_{f_{i_0}} \to
\prod\nolimits_{i_0 < i_1} A_{f_{i_0}f_{i_1}} \to
\ldots \to A_{f_1\ldots f_r}) = \colim K(A, f_1^n, \ldots, f_r^n)
$$
by More on Algebra, Lemma
\href{https://stacks.math.columbia.edu/tag/0913}{lemma extended alternating Cech is colimit koszul (Tag 0913)}.
Moreover, multiplication by $f_i^n$ on the complex
$K(A, f_1^n, \ldots, f_r^n)$ is homotopic to zero by
More on Algebra, Lemma \href{https://stacks.math.columbia.edu/tag/0663}{lemma homotopy koszul (Tag 0663)}.
Since
$$
H^q\left( LHS \right) =
\colim H^q(\text{Tot}(K^\bullet \otimes_A K(A, f_1^n, \ldots, f_r^n)))
$$
we obtain our claim. On the other hand, if $K^\bullet$ is an
object of $D_{I^\infty\text{-torsion}}(A)$, then the complexes
$K^\bullet \otimes_A A_{f_{i_0} \ldots f_{i_p}}$ have vanishing
cohomology. Hence in this case the map $LHS \to K^\bullet$
is an isomorphism in $D(A)$. The construction
$$
R\Gamma_Z(K^\bullet) =
\text{Tot}\left(
K^\bullet \otimes_A
(A \to \prod\nolimits_{i_0} A_{f_{i_0}} \to
\prod\nolimits_{i_0 < i_1} A_{f_{i_0}f_{i_1}} \to
\ldots \to A_{f_1\ldots f_r})\right)
$$
is functorial in $K^\bullet$ and defines an exact functor
$D(A) \to D_{I^\infty\text{-torsion}}(A)$ between
triangulated categories. It follows formally from the
existence of the natural transformation $R\Gamma_Z \to \text{id}$
given above and the fact that this evaluates to an isomorphism
on $K^\bullet$ in the subcategory, that $R\Gamma_Z$ is the desired
right adjoint.
\end{proof}
\end{document}
